% UTF-8 LaTeX fragment: reference answers and solutions

\sectionbanner{一、单选题}{参考答案与解析}

\begin{solution}{1}{B}{4.0}{代数与方程应用}
\analysis{把两个时刻的倍数关系同时写出，利用父子年龄差始终不变统一比较。}
\begin{enumerate}[label=\textbf{Step \arabic*:},leftmargin=*]
  \item 设儿子今年\(x\)岁，父亲今年\(y\)岁。由题意
  \[
    y-6=5(x-6),\qquad y+6=3(x+6).
  \]
  \item 两式分别化为\(y=5x-24\)与\(y=3x+12\)，所以\(2x=36\)，得\(x=18\)。
  \item 回代得父亲今年66岁。6年前为\(60=5\times12\)，6年后为\(72=3\times24\)，条件都满足。因此选B。
\end{enumerate}
\methodsummary{把跨时刻条件转成同一时刻的等式，再用不变量消元。}
\end{solution}

\begin{solution}{2}{C}{6.0}{排列组合与概率}
\analysis{按个位是否为偶数分类，再用乘法原理计数。}
\begin{enumerate}[label=\textbf{Step \arabic*:},leftmargin=*]
  \item 两位数为偶数，个位只能从\(2,4,6\)中选，有3种。
  \item 个位确定后，十位可从剩余5张卡片中任选，有5种。
  \item 不同两位数共有\(3\times5=15\)个，因此选C。
\end{enumerate}
\methodsummary{先固定受限制的位置，再安排其余位置。}
\end{solution}

\begin{solution}{3}{D}{4.0}{几何与空间}
\analysis{采用“分层剥洋葱”法。因为底面不涂色，每一层未涂色的正方体都隐藏在几何体中心。将三维立体图形逐层拆解为二维平面正方形，减去暴露在外的最外圈，剩下的即为目标数量。}

我们把这堆正方体从上往下分成5层，每层从俯视图来看，都是一个实心的正方形阵列。因为四周都暴露在外面被涂色了，我们要把每一层最外围的一圈“剥掉”。一个正方形去掉外围一圈后，它的边长左右各减去1，上下各减去1，也就是边长总共减少了2。

从上往下逐层计算内部核心未被涂色的区域：第1层全部暴露被涂色，未涂色0个；第2层全部暴露被涂色，未涂色0个；第3层边长为3，去掉外圈，内部剩下
\[
  (3-2)\times(3-2)=1
\]
个；第4层边长为4，去掉外圈，内部剩下
\[
  (4-2)\times(4-2)=4
\]
个；第5层（最底层）边长为5，四周外圈被涂色，但“底面不涂色”保护了底部。去掉外圈，剩下内部是
\[
  (5-2)\times(5-2)=9
\]
个。最后，将各层未涂色的正方体数量相加：
\[
  0+0+1+4+9=14.
\]
因此，完全未被涂色的单位正方体共有14个，正确答案选择D。

\methodsummary{寻找未涂色方块只需采用“剥洋葱法”：剥去每一层暴露在外的一圈，各层内部隐藏的方块数量即为该层原边长减2的平方，即\((n-2)^2\)，将其逐层相加即可得出总数。}
\end{solution}

\begin{solution}{4}{B}{6.0}{代数与方程应用}
\analysis{平均速度必须用总路程除以总时间，等路程两段应使用调和式而非速度算术平均。}
\begin{enumerate}[label=\textbf{Step \arabic*:},leftmargin=*]
  \item 设前后两段路程都为36千米。两段用时分别为\(36\div12=3\)小时和\(36\div18=2\)小时。
  \item 总路程为72千米，总时间为5小时。
  \item 平均速度为\(72\div5=14.4\)千米/小时，因此选B。
\end{enumerate}
\methodsummary{平均速度只由总路程与总时间决定。}
\end{solution}

\begin{solution}{5}{D}{6.0}{新定义与抽象建模}
\analysis{把7个图形看作一个周期，用序号除以周期长度后的余数定位。}
\begin{enumerate}[label=\textbf{Step \arabic*:},leftmargin=*]
  \item 图形每7个重复一次。
  \item \(2026=7\times289+3\)，所以第2026个图形与每组中的第3个相同。
  \item 每组第3个图形是\(\triangle\)，因此选D。
\end{enumerate}
\methodsummary{周期题先求余数，余数为0时才对应周期末位。}
\end{solution}

\begin{solution}{6}{E}{4.0}{几何与空间}
\analysis{利用线段平移，将内凹边向外推，把不规则图形的周长转化为最小外接矩形周长。}

题目图形由4个边长为\(5\text{ cm}\)的正方形组成。对于这类不规则的多边形，最简便的方法是使用“线段平移法”：将所有内凹的线段（横向的向两端推，纵向的向上下推）平移到外侧。这样，原图形的周长就完全等于完全覆盖它的最小外接长方形的周长。

该图形可以被切割为左、中、右三个区域。其中左侧和右侧各由1个正方形覆盖，中间区域由2个正方形竖直覆盖。宽度（横向）整体横跨了3个正方形的宽度，所以宽为
\[
  5\times3=15\text{ cm}.
\]
高度（纵向）中间区域上下堆叠了2个正方形，所以高为
\[
  5\times2=10\text{ cm}.
\]
根据长方形周长公式，代入数据得
\[
  (15+10)\times2=50\text{ cm}.
\]
因此，这个新图形的周长是\(50\text{ cm}\)，正确答案选E。

\methodsummary{利用“线段平移法”，将图形所有凹陷的边缘向外平移，不规则图形的周长问题即可转化为求其最小外接矩形的周长。}
\end{solution}

\begin{solution}{7}{C}{8.0}{数论与整除逻辑}
\analysis{正因数通常成对出现，只有完全平方数的平方根会与自身配对。}
\begin{enumerate}[label=\textbf{Step \arabic*:},leftmargin=*]
  \item 若\(m\)不是完全平方数，每个因数\(d\)都与不同的因数\(m/d\)配对，因数个数为偶数。完全平方数恰有一个自配对因数，因数个数为奇数。
  \item 不超过200的完全平方数为\(1^2,2^2,\ldots,14^2\)，共14个，因为
  \[
    14^2=196<200<225=15^2.
  \]
  \item 所求个数为\(200-14=186\)，选C。
\end{enumerate}
\methodsummary{因数配对把因数奇偶问题转成平方数计数。}
\end{solution}

\begin{solution}{8}{C}{8.0}{排列组合与概率}
\analysis{按无序数对计数，用整除条件分类或补集压缩，最后除以36个总数对。}

总的无序数对数量为
\[
  \binom92=36.
\]
乘积为6的倍数的数对为
\[
\begin{gathered}
  (1,6),(2,3),(2,6),(2,9),(3,4),(3,6),(3,8),\\
  (4,6),(4,9),(5,6),(6,7),(6,8),(6,9),(8,9),
\end{gathered}
\]
共14对。因而概率为
\[
  \frac{14}{36}=\frac7{18}.
\]
所以选C。

\methodsummary{概率计数必须让分子分母使用同一种有序或无序口径。}
\end{solution}

\begin{solution}{9}{A}{6.0}{几何与空间}
\analysis{利用“标数法”（加法原理）解题。网格中任一无障碍格子的路径数，等于到达它上方格子的路径数与左方格子的路径数之和。遇到不可穿透的障碍墙，则该方向传导的路径数视为0。}

\begin{enumerate}[label=\textbf{Step \arabic*:},leftmargin=*]
  \item \textbf{锁定唯一“关卡”。}列号单调不减，从“第2列跨到第3列”这件事在一条路径中只发生一次，记该步所在行为过闸行。墙①封第1、2行，墙②封第5、6行，故过闸行属于第3、4行。于是全局问题化为“按过闸行分类”的两类互斥且穷尽的情形。

  \item \textbf{情况一前置：过闸前（第3行过闸）。}过闸前活动域被墙①限制为左两列、上三行的\(3\times2\)小块，从\((1,1)\)走到\((3,2)\)需\(2\mathrm D+1\mathrm R\)：
  \[
    N_1^{\text{前}}=\binom31=3.
  \]

  \item \textbf{情况一后置：过闸后（扣除撞墙③者）。}过闸后在\((3,3)\)，到\((6,6)\)需\(3\mathrm D+3\mathrm R\)，不受墙限制时共
  \[
    \binom63=20.
  \]
  其中撞墙③的路径必在第3行先到达第5列（墙③在第3、4行之间的第5、6列）。从\((3,3)\)到\((3,5)\)只有\(2\mathrm R\)一种走法；再由\((3,5)\)到\((6,6)\)需\(3\mathrm D+1\mathrm R\)，共
  \[
    \binom41=4
  \]
  条，且这4条均无法下穿墙③。故合法后置数为
  \[
    N_1^{\text{后}}=\binom63-\binom41=20-4=16.
  \]

  \item \textbf{情况一总数（乘法原理）。}过闸前、后两段独立，
  \[
    N_1=N_1^{\text{前}}\times N_1^{\text{后}}=3\times16=48.
  \]

  \item \textbf{情况二前置：过闸前（第4行过闸）。}过闸前在左两列、上四行，从\((1,1)\)走到\((4,2)\)需\(3\mathrm D+1\mathrm R\)：
  \[
    N_2^{\text{前}}=\binom41=4.
  \]

  \item \textbf{情况二后置：过闸后（墙③已无关）。}过闸后已在\((4,3)\)，墙③位于第3、4行之间，已处于其下方，不可能再被挡。从\((4,3)\)到\((6,6)\)需\(2\mathrm D+3\mathrm R\)：
  \[
    N_2^{\text{后}}=\binom52=10.
  \]
  所以\(N_2=4\times10=40\)。

  \item \textbf{总数（加法原理）。}两类互斥且穷尽：
  \[
    N=N_1+N_2=48+40=\boxed{88}.
  \]
\end{enumerate}

\methodsummary{利用“标数法”（加法原理），每个格子的可达路径数等于它上方和左方相邻格子路径数之和，遇到障碍边界则对应的路径数不计入，逐格推导至终点即可。}
\end{solution}

\begin{solution}{10}{D}{8.0}{代数与方程应用}
\analysis{用局部替换证明加数只能取2或3，且三个2应改成两个3，最终按100除以3的余数构造。}

若全部加数均为1，则把其中两个1合并为2，可使乘积由1增至2，所以这种写法不可能最优。若有加数1和另一个加数\(m\ge2\)，把它们合成\(m+1\)会使乘积从\(m\)增到\(m+1\)，所以其他含1的写法也不可能最优。因此最优写法中没有加数1。

若有加数\(m\ge5\)，将其换成3与\(m-3\)，新乘积为
\[
  3(m-3)>m,
\]
所以最优写法只需2、3、4，而4可看成\(2+2\)。

三个2的乘积8小于两个3的乘积9，所以2最多出现两个。因为
\[
  100=3\times32+4,
\]
最优写法为32个3和两个2，最大乘积为
\[
  4\times3^{32}.
\]
故选D。

\methodsummary{整数和的最大乘积要用局部增益替换证明结构，再给出可达构造。}
\end{solution}

\sectionbanner{二、填空题}{参考答案与解析}

\begin{solution}{11}{19}{6.0}{排列组合与概率}
\analysis{“至少一个红球”的补事件是“两球全白”，先用补集求出概率，再按题意取最简分数的分子与分母之和。}
\begin{enumerate}[label=\textbf{Step \arabic*:},leftmargin=*]
  \item 从5个球中取2个，等可能组合总数为
  \[
    \binom52=10.
  \]
  \item 两球全白的组合数为
  \[
    \binom22=1.
  \]
  \item 至少取到1个红球的概率为
  \[
    1-\frac1{10}=\frac9{10}.
  \]
  \item 上一步所得分数已经是最简分数，所以分子与分母的和为
  \[
    9+10=19.
  \]
\end{enumerate}
\methodsummary{“至少”事件优先检查补集是否更短；求出概率后要继续完成题目指定的分数运算。}
\end{solution}

\begin{solution}{12}{3}{8.0}{数论与整除逻辑}
\analysis{从\(5!\)开始每个阶乘都含因子15，后面整段对余数没有贡献。}
\begin{enumerate}[label=\textbf{Step \arabic*:},leftmargin=*]
  \item 对任意\(m\ge5\)，\(m!\)含有因子\(3\times5=15\)，所以\(m!\)除以15余数为0。
  \item 原和与前四项之和同余：
  \[
    1!+2!+3!+4!=1+2+6+24=33.
  \]
  \item 33除以15的余数为3。
\end{enumerate}
\methodsummary{模合数时可寻找阶乘中各质因子首次齐全的位置。}
\end{solution}

\begin{solution}{13}{8}{6.0}{几何与空间}
\analysis{采用分类与分步结合的策略。首先按颜色分类（全黑或全白）。其次在同色方格中分步放置：由于不能同行同列，第一、三行的落子会相互限制，第二、四行也会独立相互限制，分别计算后再相乘。}

这是一个带有颜色和位置双重限制的排列组合问题。四个棋子颜色一致，即全在黑格或全在白格；任意两个棋子不能在同一行，也不能在同一列。

\begin{enumerate}[label=\textbf{Step \arabic*:},leftmargin=*]
  \item 先假设所有棋子都放在黑色方格中。第1行有2个黑格可供放置。一旦第1行的黑格被选中，为了满足“不同列”的要求，第3行放置棋子时，就只能选择剩下的那1个与第1行不同列的黑格。因此，在第1行和第3行放置棋子共有\(2\times1=2\)种方法。
  \item 同理，第2行的黑格与第1、3行的黑格天然不在同一列，互不干扰。第2行也有2个黑格可选；选定后，第4行只剩1个符合条件的黑格。因此，在第2行和第4行放置棋子也有\(2\times1=2\)种方法。
  \item 根据乘法原理，所有棋子都放在黑色方格中的总方法数为\(2\times2=4\)种。
  \item 黑白方格在棋盘上的分布完全对称，将所有棋子放在白色方格中同样有4种方法。根据加法原理，总方法数为
  \[
    4+4=8.
  \]
\end{enumerate}

\methodsummary{利用“分类讨论结合乘法原理”，先按黑白颜色分类，再分别计算奇偶行（或列）互不干扰的放置方法数并相乘，最后将两类结果相加即可。}
\end{solution}

\begin{solution}{14}{1012}{8.0}{数论与整除逻辑}
\analysis{分解\(2026=2\times1013\)，在1到2026中排除2或1013的倍数。}

因为
\[
  2026=2\times1013,
\]
又有
\[
  \sqrt{1013}<32.
\]
逐一试除不超过31的质数\(3,5,7,11,13,17,19,23,29,31\)，均不能整除1013，所以1013是质数。因此，与2026不互质的数恰为2的倍数或1013的倍数。

不超过2026的偶数有1013个，1013的倍数有2个，同时为二者倍数的只有2026一个。因此不互质的数有
\[
  1013+2-1=1014
\]
个。所求为
\[
  2026-1014=1012.
\]

\methodsummary{互质计数先找全部不同质因数，再对其倍数集合做容斥。}
\end{solution}

\begin{solution}{15}{174}{6.0}{几何与空间}
\analysis{本题的核心解题逻辑是体积守恒。整卷卫生纸（纯纸张部分）的体积，必定等于所有单节卫生纸展开后的长方体总体积。统一量纲：题目中交替使用\(\text{cm}\)和\(\text{mm}\)，解题第一步必须统一单位（建议统一转化为\(\text{cm}\)）。计算总体积：利用圆柱体体积公式，计算外圆柱与内部空心纸板筒的体积差，得出纯纸张的总体积。计算单节体积：将单节卫生纸看作一个薄长方体，其体积等于长\(\times\)宽\(\times\)（单层厚度\(\times\)层数）。计算总数量：用总体积除以单节纸张的体积，得到总节数，最后根据题意进行四舍五入取整。}

\begin{enumerate}[label=\textbf{Step \arabic*:},leftmargin=*]
  \item 圆柱体的体积由底面积乘以高求得。首先，统一单位为厘米：卫生纸的高度为
  \[
    100\text{ mm}=10\text{ cm}.
  \]
  整卷卫生纸的外径为\(12\text{ cm}\)，内部空心纸板筒的直径为\(5\text{ cm}\)。因此，1卷卫生纸（纯纸张部分）的体积为
  \[
    3.14\times\left(\frac{12}{2}\right)^2\times10
    -3.14\times\left(\frac{5}{2}\right)^2\times10
    =934.15\text{ cm}^3.
  \]

  \item 将单节的长和厚度转换为厘米：
  \[
    138\text{ mm}=13.8\text{ cm},\qquad
    0.13\text{ mm}=0.013\text{ cm}.
  \]
  因为每节卫生纸有3层，所以每一节的体积为
  \[
    10\times13.8\times0.013\times3=5.382\text{ cm}^3.
  \]

  \item 将总纸张体积除以单节体积：
  \[
    \frac{934.15}{5.382}\approx173.57.
  \]
  根据题目要求四舍五入到整数，即为174节。
\end{enumerate}

\methodsummary{把握“体积守恒”核心，在统一单位的前提下，用空心圆柱纸张的总净体积除以单节展开后的长方体体积，即可求出总节数。}
\end{solution}

\sectionbanner{三、综合题}{参考答案与解析}

\begin{solution}{16（1）}{B}{4.0}{新定义与抽象建模}
\analysis{紧扣古埃及分数的“两条铁律”进行排查。铁律1要求分子必须全为1；铁律2要求分母绝对不能重复。只需逐一检验选项是否违规即可得出唯一正确答案。}

A选项：
\[
  \frac34=\frac14+\frac14+\frac14,
\]
分母重复了三个4，违反铁律2，排除。

B选项：
\[
  \frac34=\frac12+\frac14,
\]
分子全为1，分母分别是2和4，没有重复，完全符合两条铁律，正确。

C选项：
\[
  \frac34=\frac24+\frac14,
\]
其中\(\frac24\)分子不是1，违反铁律1，排除。

D选项：
\[
  \frac34=\frac13+\frac13+\frac1{12},
\]
分母重复了两个3，违反铁律2，排除。

\methodsummary{严格对照“分子必须为1且分母不重复”的两条新定义铁律，逐一“排雷”筛选选项。}
\end{solution}

\begin{solution}{16（2）}{42}{6.0}{新定义与抽象建模}
\analysis{第一个空填7；第二个空填42。这是对“魔法拆分公式”的直接运用。明确公式里的变量\(n=6\)，代入公式计算出两个新分母即可，注意题目要求数字按从小到大填空。}

已知魔法拆分公式：
\[
  \frac1n=\frac1{n+1}+\frac1{n\times(n+1)}.
\]
题目要拆分\(\frac16\)，令\(n=6\)。代入公式得分母分别为
\[
  6+1=7,\qquad 6\times7=42.
\]
根据题目要求“第一个比第二个小”，所以两空分别填7和42，第二个括号里填写42。

\methodsummary{运用“代入法”，将具体数字直接替换公式中的字母\(n\)，按运算顺序精准计算出新分母。}
\end{solution}

\begin{solution}{16（3）}{30}{6.0}{新定义与抽象建模}
\analysis{三个空依次填入5、6、30。核心策略是“先分后变”。先强行拆成两个\(\frac15\)相加，再利用公式对其中一个\(\frac15\)进行变身，最后按数字从小到大排序填空。}

\begin{enumerate}[label=\textbf{Step \arabic*:},leftmargin=*]
  \item 先强行拆分：
  \[
    \frac25=\frac15+\frac15.
  \]
  此时分母重复。
  \item 保留第一个\(\frac15\)不动，对第二个使用公式（令\(n=5\)），得
  \[
    \frac15=\frac16+\frac1{30}.
  \]
  \item 合并替换得
  \[
    \frac25=\frac15+\frac16+\frac1{30}.
  \]
  按从小到大顺序，分母依次为5、6、30。
\end{enumerate}

\methodsummary{采用“先分后变”的转化策略，先同分母强行拆分，再套用公式对局部进行等量替换。}
\end{solution}

\begin{solution}{17（1）}{C}{4.0}{新定义与抽象建模}
\analysis{将边\(CD\)视为透明窗户，利用镜面反射原理。寻找起点\(A\)关于\(CD\)的镜像对称点\(A'\)，把反弹轨迹等效为从\(A'\)发出的直射光束，扫射边界即可。}

找镜像点：在正六边形中，对角线\(AC\)与边\(CD\)恰好垂直，因为
\[
  120^\circ-30^\circ=90^\circ.
\]
若把\(CD\)当作镜面向下翻折，\(A\)的镜像点\(A'\)正好位于\(A\)到\(C\)连线的延长线上。

光束扫射：从\(A'\)点向“窗户”\(CD\)打出光束。当光线擦着极左端\(C\)点射出时，沿直线\(A'C\)会笔直命中顶点\(A\)。当光线擦着极右端\(D\)点射出时，沿直线\(A'D\)连线延长，恰好精准命中顶点\(E\)。

反弹后的光束呈扇形，其辐射范围被严密限制在顶点\(A\)到顶点\(E\)之间。因此，这束光只能打在正六边形右侧的边\(AF\)和边\(EF\)上。

\methodsummary{遇到折线反射，把反射面当透明窗户，从起点的“镜像分身”拉直连线，寻找光束两端的极值边界。}
\end{solution}

\begin{solution}{17（2）}{B}{6.0}{新定义与抽象建模}
\analysis{2次反弹相当于跨越2层镜像。将终点\(D\)往下翻折两次得到终极目标\(D_2\)。连接起点\(A\)与\(D_2\)，形成拉直的轨迹，利用相似三角形的比例求解。}

\begin{enumerate}[label=\textbf{Step \arabic*:},leftmargin=*]
  \item \textbf{建构楼层。}设正六边形的平行对边间距为1个“楼层高”。起点\(A\)位于第一面镜子\(CD\)上方1层。
  \item \textbf{定位终点。}终点\(D\)经过2次镜像翻折（向下延展），其最终的镜像点\(D_2\)位于\(CD\)镜面下方2层的位置。
  \item \textbf{拉直连线。}将真实的曲折轨迹拉直为直线\(AD_2\)，该直线穿过\(CD\)边交于点\(X\)。
  \item \textbf{比例分配。}直线在垂直方向上的落差比为\(1:2\)（上方1层、下方2层）。根据相似三角形等比例放大原理，它在水平方向上分得的距离比例也为\(1:2\)。
  \item \textbf{计算结果。}边\(CD\)的总长度为\(1\text{ cm}\)，因此\(CX\)占水平总长度的
  \[
    \frac1{1+2},
  \]
  即
  \[
    CX=\frac13\text{ cm}.
  \]
\end{enumerate}

\methodsummary{将多次反弹转化为多层镜面翻折，连接起点与最终镜像点，用“垂直楼层数之比＝水平路程之比”求解。}
\end{solution}

\begin{solution}{17（3）}{C}{6.0}{新定义与抽象建模}
\analysis{3次反弹需翻折3层找到目标\(D_3\)。在密铺网格中精准数出\(D_3\)的“水平总位移”和“垂直总落差”，再次利用相似三角形的比例分配求解。}

\begin{enumerate}[label=\textbf{Step \arabic*:},leftmargin=*]
  \item \textbf{寻找第三层终点。}经过3次镜面翻折，终点\(D\)被映射到了更远更深的第3层镜像点\(D_3\)。
  \item \textbf{测算总跨度。}借助六边形密铺网格观察，目标\(D_3\)相对于参考点\(C\)，在水平方向上向右偏移了2.5个边长，即水平总宽度为\(2.5\text{ cm}\)。
  \item \textbf{测算总落差。}在垂直方向上，起点\(A\)在\(CD\)上方1层，而\(D_3\)位于\(CD\)下方2.5层，因此总垂直落差为
  \[
    1+2.5=3.5
  \]
  层。
  \item \textbf{相似比例。}直线\(AD_3\)穿过\(CD\)边交于点\(Y\)。直线在通过第一面镜子前，走过的垂直比例是
  \[
    \frac1{3.5}=\frac27.
  \]
  \item \textbf{计算结果。}水平方向的路程\(CY\)也占总水平宽度的\(\frac27\)，即
  \[
    CY=2.5\times\frac27=\frac57\text{ cm}.
  \]
\end{enumerate}

\methodsummary{无论反弹多少次，核心就是锁定最终镜像点的“总水平宽度”和“总垂直落差”，用单一比例算式求解截距。}
\end{solution}

\begin{solution}{18（1）}{C}{4.0}{新定义与抽象建模}
\analysis{到达某格的走法数，等于到达它前面一、二、三格的走法数之和。依次写出前几个格子的走法即可求出。}

设到达第\(n\)格的走法为\(W(n)\)。
\begin{enumerate}[label=\textbf{Step \arabic*:},leftmargin=*]
  \item 走1步到第1格，只有走1一种，即\(W(1)=1\)。
  \item 到第2格，可走\(1+1\)或直接跨2，即\(W(2)=2\)。
  \item 到第3格，可从起点跨3，或从第1、第2格跨来，即
  \[
    W(3)=1+1+2=4.
  \]
  \item 到第4格，只能从第1、2、3格跨来，故
  \[
    W(4)=W(1)+W(2)+W(3)=1+2+4=7.
  \]
\end{enumerate}
故有7种走法。

\methodsummary{运用递推法（或动态规划思想），提炼出状态转移方程，即“到达当前格的走法数等于到达它前面三格的走法数之和”，并穷举出前三格的初始基础走法。}
\end{solution}

\begin{solution}{18（2）}{274}{6.0}{新定义与抽象建模}
\analysis{基于第1问的递推关系，连续将前三项的数值相加得到下一项，一步步向后推导直至第10项。}

由上问知核心递推公式为
\[
  W(n)=W(n-1)+W(n-2)+W(n-3).
\]
已知前四项为
\[
  W(1)=1,\quad W(2)=2,\quad W(3)=4,\quad W(4)=7.
\]
依次推导后序项：
\[
\begin{aligned}
  W(5)&=2+4+7=13,\\
  W(6)&=4+7+13=24,\\
  W(7)&=7+13+24=44,\\
  W(8)&=13+24+44=81,\\
  W(9)&=24+44+81=149,\\
  W(10)&=44+81+149=274.
\end{aligned}
\]
所以，到达第10格共有274种走法。

\methodsummary{运用递推法（或动态规划思想），提炼出状态转移方程，即“到达当前格的走法数等于到达它前面三格的走法数之和”，并穷举出前三格的初始基础走法。}
\end{solution}

\begin{solution}{18（3）}{168}{6.0}{新定义与抽象建模}
\analysis{将路线拆分为“起点到第6格”和“第6格到第10格”两段独立的子路径。算出各自走法数后相乘。}

\begin{enumerate}[label=\textbf{Step \arabic*:},leftmargin=*]
  \item 第一阶段，从起点走到第6格。根据上问推导的递推结果，走法有\(W(6)=24\)种。
  \item 第二阶段，从第6格走到第10格。这段路程的跨度为
  \[
    10-6=4
  \]
  格，这与从起点出发走4格完全等价，走法有\(W(4)=7\)种。
  \item 由于这两个阶段连续发生且相互独立，根据分步乘法原理，总走法数为
  \[
    W(6)\times W(4)=24\times7=168.
  \]
\end{enumerate}
故共有168种符合条件的走法。

\methodsummary{运用分步乘法原理。将整个运动轨迹在“第6格”处切断，拆分为两个独立阶段。先求出从起点到达第6格的走法数，再求出从第6格到达第10格（相当于从头走4格）的走法数，最后将两个阶段的结果相乘即可。}
\end{solution}
